% Grado 10 -- generado por tools/build_tex.py a partir de raw/grado_10.txt
\section{Grado 10}\label{grado:10}

% PDF p. 34
\dba{1}{Comprende, que el reposo o el movimiento rectilíneo uniforme, se presentan cuando las fuerzas aplicadas sobre el sistema se anulan entre ellas, y que en presencia de fuerzas resultantes no nulas se producen cambios de velocidad.}

\evidencias

\begin{itemize}
  \item Predice el equilibrio (de reposo o movimiento uniforme en línea recta) de un cuerpo a partir del análisis de las fuerzas que actúan sobre él (primera ley de Newton).
  \item Estima, a partir de las expresiones matemáticas, los cambios de velocidad (aceleración) que experimenta un cuerpo a partir de la relación entre fuerza y masa (segunda ley de Newton).
  \item Identifica, en diferentes situaciones de interacción entre cuerpos (de forma directa y a distancia), la fuerza de acción y la de reacción e indica sus valores y direcciones (tercera ley de Newton).
\end{itemize}

\ejemplo

\fig{g10_p34_1}{Ejemplo (solo imagen): un adulto y un niño tiran de una misma cuerda en sentidos opuestos; el adulto está sobre una superficie resbaladiza (hielo o agua). Situación para analizar las fuerzas, el equilibrio y el par acción-reacción (leyes de Newton).}


% PDF p. 34
\dba{2}{Comprende la conservación de la energía mecánica como un principio que permite cuantificar y explicar diferentes fenómenos mecánicos: choques entre cuerpos, movimiento pendular, caída libre, deformación de un sistema masa-resorte.}

\evidencias

\begin{itemize}
  \item Predice cualitativa y cuantitativamente el movimiento de un cuerpo al hacer uso del principio de conservación de la energía mecánica en diferentes situaciones físicas.
  \item Identifica, en sistemas no conservativos (fricción, choques no elásticos, deformación, vibraciones) las transformaciones de energía que se producen en concordancia con la conservación de la energía.
\end{itemize}

\ejemplo

\fig{g10_p34_2}{Ejemplo (solo imagen): dos jóvenes juegan billar; el taco golpea la bola blanca, que choca con otra bola. Situación de choque entre cuerpos para analizar la conservación y las transformaciones de la energía.}


